\section{Background and Related Work}
\label{sec:background}

\subsection{What the WireGuard handshake computes}

WireGuard's handshake is Noise \texttt{IKpsk2}, instantiated with X25519,
ChaCha20-Poly1305 and BLAKE2s. The initiator knows the
responder's static public value in advance---this is what makes the pattern
\texttt{IK} rather than \texttt{XX}---and the exchange completes in two
messages, an initiation of \wgInitBytes\,B and a response of \wgRespBytes\,B
(Figure~\ref{fig:flow}).

\begin{figure}[t]
\centering
\resizebox{\columnwidth}{!}{\begin{tikzpicture}[x=1cm,y=1cm,font=\footnotesize]
  \draw[black!25,thick] (0,-0.25) -- (0,-4.5);
  \draw[black!25,thick] (6.2,-0.25) -- (6.2,-4.5);
  \node[anchor=south] at (0,-0.2) {\bfseries Initiator};
  \node[anchor=south] at (6.2,-0.2) {\bfseries Responder};
  \node[anchor=north,black!55,font=\scriptsize] at (0,-0.3)
    {static $(s_i,S_i)$, psk};
  \node[anchor=north,black!55,font=\scriptsize] at (6.2,-0.3)
    {static $(s_r,S_r)$, psk};

  \draw[->,thick] (0,-1.35) -- (6.2,-1.35);
  \node[anchor=south,inner sep=2pt] at (3.1,-1.32)
    {\bfseries handshake initiation, 148\,B};
  \node[anchor=north,black!55,font=\scriptsize,align=center,inner sep=2pt]
    at (3.1,-1.4)
    {$E_i$ $\mid$ enc\_static (48\,B) $\mid$ enc\_timestamp (28\,B)\\
      $\mid$ mac1 $\mid$ mac2};
  \node[anchor=west,black!55,font=\scriptsize] at (0.1,-2.25)
    {$es=\mathrm{DH}(e_i,S_r)$,\quad $ss=\mathrm{DH}(s_i,S_r)$};

  \draw[<-,thick] (0,-3.0) -- (6.2,-3.0);
  \node[anchor=south,inner sep=2pt] at (3.1,-2.97)
    {\bfseries handshake response, 92\,B};
  \node[anchor=north,black!55,font=\scriptsize,inner sep=2pt] at (3.1,-3.05)
    {$E_r$ $\mid$ enc\_nothing (16\,B) $\mid$ mac1 $\mid$ mac2};
  \node[anchor=east,black!55,font=\scriptsize] at (6.1,-3.62)
    {$ee=\mathrm{DH}(e_r,E_i)$,\quad $se=\mathrm{DH}(e_r,S_i)$};
  \node[anchor=east,okgreen,font=\scriptsize] at (6.1,-3.92)
    {psk mixed by $\mathrm{KDF}_3$};

  \draw[okgreen,dashed,thick] (0,-4.35) -- (6.2,-4.35);
  \node[anchor=south,okgreen,font=\scriptsize,inner sep=2pt] at (3.1,-4.32)
    {transport keys $(T_i,T_r)=\mathrm{KDF}_2(ck,\varepsilon)$};
\end{tikzpicture}
}
\caption{The two handshake messages, their sizes on the wire, and the
Diffie-Hellman terms each contributes to the key schedule. Sizes are computed
from the field layout of Section~\ref{sec:modelling}, not transcribed.}
\label{fig:flow}
\end{figure}

Four Diffie-Hellman terms enter the key schedule, and every property in this
paper turns on which of them an attacker can compute:

\begin{itemize}\itemsep2pt
\item \(es=\mathrm{DH}(e_i,S_r)\), initiator ephemeral against responder
      static, binding message~1 to its intended recipient;
\item \(ss=\mathrm{DH}(s_i,S_r)\), static-static, which gives \texttt{IK} its
      first-message authentication;
\item \(ee=\mathrm{DH}(e_r,E_i)\), ephemeral-ephemeral, the term that carries
      forward secrecy;
\item \(se=\mathrm{DH}(e_r,S_i)\), responder ephemeral against initiator
      static, which authenticates the initiator to the responder in
      message~2.
\end{itemize}

A pre-shared key, when configured, is mixed in by \(\mathrm{KDF}_3\) during
message~2, yielding a chaining key, a value \(\tau\) folded into the
transcript hash, and an AEAD key. Because \(\tau\) enters the transcript, a
mismatched pre-shared key makes the authenticated decryption in message~2 fail
rather than producing a silently divergent session.

\para{Beyond the Noise pattern.} WireGuard adds a TAI64N timestamp, session
indices, and two MAC fields. The first, \texttt{mac1}, is keyed by a hash of
the recipient's static public value; it demonstrates knowledge of that value
and suppresses unsolicited traffic, but is not a cryptographic authenticator.
The second, \texttt{mac2}, carries a cookie under load. Both are
consequential for the correspondence of Section~\ref{sec:modelling} and for
what the symbolic model can and cannot say.

\subsection{Symbolic verification, and why two provers}

Both tools used here work in the Dolev-Yao model~\cite{dolev1983security}: the
attacker controls the network completely and the cryptographic primitives are
perfect.

\para{ProVerif.} ProVerif~\cite{blanchet2016proverif} translates an applied
pi-calculus process into Horn clauses and saturates them. It is fully
automatic and fast, handles an unbounded number of sessions, and
over-approximates: a proof is sound, while a reported attack may occasionally
be spurious. Saturation either converges or it does not, and when it does not
the tool gives no partial answer.

\para{Tamarin.} Tamarin~\cite{meier2013tamarin} works with multiset rewriting
and constraint solving, represents mutable state directly, and normalises
Diffie-Hellman terms natively. It reasons backwards from the property rather
than forwards from the rules, which turns out to matter for exactly the
queries ProVerif could not settle.

\para{Stating authentication precisely.} Authentication is stated throughout
in Lowe's hierarchy~\cite{lowe1997hierarchy}: aliveness, weak agreement,
non-injective agreement, injective agreement. Naming the level explicitly
matters here, because the key-compromise result of
Section~\ref{sec:results} is precisely a case of two adjacent levels
separating.

\subsection{Related work}

\para{Prior analyses of WireGuard.} Donenfeld and
Milner~\cite{donenfeld2018formal} give a symbolic analysis in Tamarin covering
the handshake and its state machine. Lipp, Blanchet and
Bhargavan~\cite{lipp2019cryptoverif} give a mechanised computational proof in
CryptoVerif, which yields concrete bounds this work cannot. Our contribution
relative to both is not a stronger result on the same property, but the
deliberate pairing of two symbolic tools on one transcription, and the
extension of the analysis past the cryptography into path reachability.

\para{Post-quantum WireGuard.} Hülsing et
al.~\cite{hulsing2021pqwireguard} redesign the handshake around KEMs directly,
and Rosenpass~\cite{rosenpass2023} instead runs a post-quantum key agreement
on a separate channel and injects the result into WireGuard's pre-shared key
slot. Section~\ref{sec:pq} does not propose a third design; it supplies the
arithmetic that explains why the second shape exists.

\para{KEM binding.} Cremers, Dax and Medinger~\cite{cremers2024kembinding}
systematise binding properties for KEMs and provide the vocabulary in which
the result of Section~\ref{sec:results} is stated. The property that separates
our two KEM models is MAL-BIND-K-PK.
